% !TEX root = ../main.tex

% --- Résumé (Langue principale) ---
\chapter*{Résumé}
% Décommente la ligne suivante si tu veux que le résumé apparaisse dans la table des matières :
% \addcontentsline{toc}{chapter}{Résumé}

Ici, tu peux insérer le résumé de ton mémoire. Il doit synthétiser ta problématique, ta méthodologie, tes principaux résultats et tes conclusions en un ou deux paragraphes. C'est souvent la dernière chose que l'on rédige !

\vspace{1cm}
\noindent \textbf{Mots-clés :} Mot-clé 1, Mot-clé 2, Mot-clé 3, Mot-clé 4.

\vspace{2cm}

% --- Abstract (Version anglaise, optionnelle mais recommandée) ---
\begin{center}
    \Large \textbf{Abstract}
\end{center}
\vspace{0.5cm}

Here you can insert the English translation of your abstract. 

\vspace{1cm}
\noindent \textbf{Keywords:} Keyword 1, Keyword 2, Keyword 3, Keyword 4.